\section{Conclusion}\label{sec:conclusion}

If local summaries preserve the evidence a task oracle needs, and local merges preserve the cross-span interactions, then preference optimization can be carried out on compression trees rather than full documents. The theorem ladder makes this precise: the structural bundle gives root preservation, the optimization bundle gives objective equivalence for DPO/GRPO, and the certificate bundle turns sampled audits into a finite-sample bound on preference drift. Beyond preservation, the same local signal the audit samples is a signal the model can be supervised on, so when a task's structure aligns with the local laws, local supervision becomes a first-class resource alongside root labels. The Markov budget-split experiments make this concrete: when the oracle decomposes into additive leaf-local terms plus a boundary correction, roughly half the root labels can be replaced by local labels at equal total cost without loss of root accuracy.

The paper arcs from the classical to the learned case along a single gradient. Section~\ref{sec:mergeable-sketches} grounds the framework in the data-systems literature on mergeable sketches. Sections~\ref{sec:markov-interlude}--\ref{sec:markov-walkthrough} show a task whose oracle defeats any commutative sketch and the tree-plus-FNO architecture that resolves it. Section~\ref{sec:hll-parity} closes the classical-limit loop empirically: running the framework with a deterministic register-wise merge recovers HyperLogLog byte-for-byte on a register-based implementation, and a fully end-to-end learned version tracks the classical sketch within a small multiple of its theoretical relative standard error. Section~\ref{sec:manifesto-llm} then takes the same machinery to Manifesto-RILE scoring, where no hand-designed sketch exists and the sufficient statistic is a high-dimensional learned representation. The Semantic Forests companion extends these ideas to orchestration and deployment at system scale.
